% Modernised reading — fonds Alexandre Grothendieck, shelfmark 31, pages 1-100.
%
% This is an INTERPRETATION, not a transcription. It re-reads the pages in
% today's notation and vocabulary, states the hypotheses the manuscript leaves
% implicit, and drops the critical apparatus so that the mathematics reads
% straight through. Everything the brackets and underlines carried — a doubtful
% reading, a notation he used differently, a missing passage — moves into a
% footnote.
%
% It is held to being mathematically correct as it stands, not merely faithful
% to the page. Where the manuscript is loose or mistaken, the true statement is
% what appears here, and the footnote says what the page has. In any dispute of
% substance the transcription governs: whoever cites Grothendieck must cite the
% batch-NN.fr.tex files.
%
% Scope: the folder was read WHOLE before any of this was written — all five
% batches, pages 1-100, in one pass — and it is written as ONE file for the
% whole shelfmark. The folder's sixth batch is a single leaf, page 101, which
% continues the foreign typescript of page 99 (Bourbaki, Algèbre commutative,
% chap. VII, n° 339, on the essential valuations of a Krull ring) and carries
% nothing in his hand; it has no transcription and nothing here reads it.
%
% The spine: the folder is the workshop of ONE question — how high can the
% cohomology of a scheme climb, and in particular does it stop at the dimension
% for an affine scheme. That is Artin's affine vanishing theorem, and page 76
% (Lemme 3) is where the folder reaches it. Everything else is either an input
% to that (the cohomological dimension of fields, pages 2-8, 19-22, 46-51), a
% tool for it (Čech against derived functors, 24-30; the trace, 41-44; the
% negative answer about injective sheaves, 54-61), a refinement of it (the
% function phi, 64-70; the two induction rungs, 73-80; the conjectures, 81-84),
% or a consequence drawn from it (Lefschetz, 89-90; local asphericity, 92-98).
%
% One notation is fixed here for the whole folder and held to throughout, in
% the "Conventions" section: the folder's E^d_k and Y[t_1,…,t_r] are written
% A^d_k and A^r_Y, and « préschéma » is written « schéma », since the manuscript
% predates the change of vocabulary in EGA.
%
% Four departures from the pages are substantive rather than cosmetic, and each
% is footnoted where it occurs:
%   - page 73's Théorème 1 is stated for « X = Spec A » with A strictly local,
%     which would make it vacuous (a strictly local scheme has no higher
%     cohomology at all); the reading says so and does not repair the statement,
%     which the leaf does not determine;
%   - page 50's Corollaire 2 is stated for a finite group with trivial action;
%     the argument gives it for every finite Galois module, and the reading
%     states the stronger true form;
%   - page 28's conditions (ii) and (iii) carry a range of indices overwritten
%     on the leaf and doubtful; the reading gives the range the proof uses and
%     footnotes the ambiguity rather than choosing silently;
%   - page 6 writes H^i(k, prod_{k'/k} G') ≃ H^i(k,G) where Shapiro's lemma
%     requires H^i(k',G'); the reading writes the latter.
%
% Pass: Opus 5 (claude-opus-5), 2026-09-11 — first pass, unchecked against the
% pages by a human.
\documentclass[11pt,a4paper]{article}
% Le préambule commun à toutes les transcriptions du fonds Grothendieck.
%
% Il tient en une idée : **l'incertitude se compose, elle ne se résout pas.**
% Une transcription qui lisse un mot douteux a détruit la seule chose pour
% laquelle on la faisait — un lecteur doit pouvoir savoir, à chaque endroit,
% ce qui était sur le feuillet et ce qui a été ajouté. Les six macros qui
% suivent sont donc l'appareil critique, et non de la décoration.
%
% Les mêmes commandes sont reconnues par `scripts/render.mjs`, qui produit la
% vue de lecture du site. Une macro ajoutée ici doit l'être aussi là-bas,
% faute de quoi le rendu échouera bruyamment — ce qui est voulu : un rendu
% silencieusement faux est pire qu'un rendu absent.

\ProvidesPackage{grothendieck}[2026/08/08 Transcriptions du fonds Grothendieck]

\usepackage{amsmath,amssymb,amsthm}
\usepackage{mathrsfs}
\usepackage{stmaryrd}
% Les diagrammes commutatifs sont partout dans ce fonds : c'est la langue
% même de la géométrie algébrique de Grothendieck, et une transcription qui
% les rendrait en prose ne transcrirait rien.
\usepackage{tikz-cd}
% Français par défaut — c'est la langue des feuillets — mais l'anglais est
% chargé aussi : la traduction et le résumé sont des documents anglais, et la
% typographie française (espaces avant les deux-points) y serait une faute.
% Ces documents basculent par \selectlanguage{english} en tête de corps.
\usepackage[english,french]{babel}
\usepackage{fontspec}
\usepackage{geometry}
\geometry{a4paper,margin=25mm,marginparwidth=28mm}
\usepackage{marginnote}
\usepackage{xcolor}
% ulem, not soul. `soul`'s \st and \ul are built on a character-by-character
% scan that XeLaTeX's font machinery breaks: the document fails with an
% undefined control sequence rather than degrading. ulem does the same two jobs
% and is Unicode-engine safe. `normalem` stops it hijacking \emph, which the
% transcriptions use for Grothendieck's own emphasis.
\usepackage[normalem]{ulem}
\usepackage{hyperref}
\hypersetup{colorlinks=true,linkcolor=black,urlcolor=black,pdfborder={0 0 0}}

% --- Métadonnées du lot -----------------------------------------------------
% Reprises telles quelles par le script de rendu, qui les affiche en tête de
% la vue de lecture. Elles fixent ce que le fichier prétend couvrir : sans
% elles, une transcription flotte sans point d'attache dans un fonds de seize
% mille feuillets.

\newcommand{\folder}[1]{\gdef\grothFolder{#1}}
\newcommand{\batch}[1]{\gdef\grothBatch{#1}}
\newcommand{\leaves}[2]{\gdef\grothFirst{#1}\gdef\grothLast{#2}}
\newcommand{\foldertitle}[1]{\gdef\grothFoldertitle{#1}}
\gdef\grothFolder{?}\gdef\grothBatch{?}\gdef\grothFirst{?}\gdef\grothLast{?}\gdef\grothFoldertitle{}

% --- Le résumé --------------------------------------------------------------
%
% Il ouvre la lecture modernisée, sous le titre « Résumé », et il est composé
% en petit. Ce n'est pas un ornement : c'est le seul passage du document qui
% ne soit pas des mathématiques, et un lecteur doit pouvoir le reconnaître
% avant de l'avoir lu — puis le sauter s'il connaît déjà le sujet, ou s'y
% arrêter s'il ne le connaît pas. Le filet qui le suit marque la reprise du
% corps.

\newenvironment{resume}
  {\small\setlength{\parskip}{0.45em}}
  {\par\medskip\hrule\medskip}

% --- Mots-clés ---------------------------------------------------------------
%
% En anglais, en clôture du résumé de la lecture modernisée : le vocabulaire
% moderne sous lequel chercher ce que les feuillets construisent. Le manifeste
% du site (`npm run manifest`) les extrait pour étiqueter la cote entière —
% cette ligne est la source unique des tags, il n'y a pas de fichier à côté.
\newcommand{\keywords}[1]{\par\smallskip\noindent
  {\footnotesize\textsc{Keywords} --- \textit{#1}}\par}

% --- L'appareil critique ----------------------------------------------------

% Le numéro de feuillet, en marge. C'est l'ancre qui permet de régler un
% désaccord sur une formule en regardant *un* feuillet plutôt qu'une chemise
% de six cents. Le site s'en sert aussi pour tourner les pages du fac-similé
% quand on fait défiler la transcription.
\newcommand{\leaf}[1]{%
  \marginnote{\footnotesize\color{gray!70}\textsf{#1}}%
  \label{leaf:#1}\ignorespaces}

% Illisible : on ne devine pas. Un mot inventé qui se lit comme les autres est
% le pire résultat possible.
\newcommand{\ill}{\textcolor{red!60!black}{[\,\ldots\,]}}

% Lecture douteuse : le mot est donné, et signalé comme incertain.
\newcommand{\uncertain}[1]{\uline{#1}}

% Ajout de l'éditeur — une lettre manquante, un mot sous-entendu.
\newcommand{\add}[1]{\textcolor{blue!50!black}{[#1]}}

% Biffé par Grothendieck lui-même. Ce qu'il a rayé fait partie du document :
% une rature dit souvent où la pensée a changé de direction.
\newcommand{\struck}[1]{\sout{#1}}

% Note du transcripteur, sur l'état matériel du feuillet.
\newcommand{\note}[1]{\par\smallskip{\small\color{gray!60!black}\textit{[#1]}}\par\smallskip}

% Note marginale de Grothendieck, distincte du corps du texte.
\newcommand{\marginal}[1]{\par\smallskip\noindent\hspace{1em}%
  {\small\color{gray!60!black}#1}\par\smallskip}

% --- Titre ------------------------------------------------------------------

\AtBeginDocument{%
  \begin{center}
    {\large\bfseries Fonds Alexandre Grothendieck --- cote n\textsuperscript{o}~\grothFolder}\\[2pt]
    {\itshape\grothFoldertitle}\\[4pt]
    {\small feuillets \grothFirst--\grothLast\ (lot \grothBatch)}\\[2pt]
    {\footnotesize\color{gray!60!black}Transcription établie d'après le fac-similé numérisé
     par l'Université de Montpellier.\\
     Les lectures incertaines sont soulignées, les illisibles notées [\,\ldots\,],
     les ajouts entre crochets.}
  \end{center}
  \vspace{1em}}

\endinput


\folder{31}
\pages{1}{100}
\dating{[1963-1973]}
\watermark{Édition de démonstration\\interprétation personnelle de l'œuvre}
\foldertitle{SGA A [SGA 4] (Résidus de rédaction) : notes manuscrites (s.d.) — lecture modernisée du dossier entier}

\begin{document}

\section*{Résumé}

\begin{resume}
La chemise porte, de la main des archivistes, « SGA A [SGA 4] (Résidus de
rédaction) ». Résidus : ce qui reste sur la table quand un traité est parti
chez le dactylographe. Cent pages, écrites au dos de lettres et de tapuscrits,
qui n'ont pas l'air d'un texte et n'en sont pas un. Mais une seule question les
traverse, et il vaut la peine de l'avoir en tête avant d'en ouvrir une.

Voici cette question. À un espace, on sait associer une suite de groupes
$H^0, H^1, H^2, \dots$ qui comptent ses trous : $H^0$ ses morceaux, $H^1$ ses
lacets qui ne se remplissent pas, $H^2$ ses cavités, et ainsi de suite. Pour un
espace ordinaire cette suite s'arrête, et elle s'arrête au rang qui mesure la
taille de l'espace : une surface n'a pas de $H^3$, il n'y a pas de place pour
une cavité de dimension trois dans un objet qui n'en a que deux. C'est une
évidence tant qu'on travaille avec des espaces ordinaires, et l'on cesse d'y
penser.

Or vers 1960 on venait de construire, pour les objets de la géométrie
algébrique, une cohomologie neuve — la cohomologie étale — faite pour compter
des trous que la topologie de Zariski ne voyait pas. Neuve, elle n'avait
aucune des évidences des anciennes. Rien ne disait où sa suite s'arrêtait, ni
même qu'elle s'arrêtât. Et tout en dépendait : les théorèmes de finitude, les
suites spectrales, les conjectures de Weil que la théorie était faite pour
attaquer, tous demandent un rang au-delà duquel il n'y a plus rien. Le nombre
qui répond s'appelle la \emph{dimension cohomologique}, et ce dossier est
l'atelier où on le cerne.

L'idée qui arrive tient dans une remarque sur les espaces ordinaires. Un plan
tout entier n'a aucun trou : il se rétracte sur un point. Une sphère en a un,
mais il suffit de lui ôter un point — de la rendre \emph{ouverte} — pour qu'il
disparaisse. En géométrie algébrique, les objets « ouverts » de cette sorte
s'appellent les schémas affines, et l'analogie suggère qu'ils devraient être
cohomologiquement pauvres : non pas sans trous, mais avec beaucoup moins que
leur taille n'autoriserait. Le pari du dossier est précis : \emph{un schéma
affine de dimension $n$ sur un corps séparablement clos n'a pas de cohomologie
au-delà du rang $n$}. C'est le théorème d'annulation d'Artin, et la page 76 est
l'endroit où le dossier y arrive.

La stratégie pour l'atteindre est celle dont Grothendieck a fait une méthode.
On ne prouve pas le théorème sur l'objet : on le fait descendre le long d'une
fibration jusqu'à des objets sans épaisseur. Un espace affine de dimension $d$
se projette sur un espace affine de dimension $d-1$, les fibres sont des
droites ; si l'on sait ce que vaut la cohomologie des fibres, une suite
spectrale recolle et l'on a gagné un cran. Descendue jusqu'au bout, la question
n'est plus sur des schémas mais sur des \emph{corps} — le plus simple des
objets géométriques, un point sans épaisseur — et c'est là que le dossier passe
le plus clair de son temps. Combien de cohomologie a un corps ? La réponse
dépend de son arithmétique et non de sa géométrie, et elle est tantôt facile
(un corps de fonctions d'une variable sur un corps algébriquement clos en a
très peu), tantôt féroce : un corps où $-1$ n'est pas une somme de carrés — les
nombres réels, par exemple — en a infiniment, à tous les rangs, pour toujours.
Cette dernière pathologie occupe deux pages du dossier, et elle explique
pourquoi tant d'énoncés portent « $\ell$ premier à la caractéristique » ou « $k$
séparablement clos » : ce ne sont pas des précautions d'écriture, ce sont les
endroits où la théorie casse.

Le dossier n'est pas une démonstration ; c'est un chantier, et il faut le lire
comme tel. Des énoncés y sont écrits trois fois de suite, chaque version
raturant la précédente. Des preuves s'arrêtent au milieu d'une phrase. Une
question posée dans une marge — « cette restriction est-elle essentielle
??? » — reste sans réponse trente pages plus loin. C'est peut-être ce que ces
feuillets ont de plus rare : non pas les théorèmes, qu'on trouvera rédigés
ailleurs, mais la trace de ce que c'est que de ne pas encore savoir, et de le
noter dans la marge plutôt que de passer outre.

Pour aller plus loin, les noms d'aujourd'hui : théorème d'annulation d'Artin,
dimension cohomologique d'un corps, théorème de Tsen, corps $C_1$, cohomologie
de Tate, lemme d'Abhyankar, acyclicité locale des morphismes lisses.

\keywords{cohomological dimension, étale cohomology, Artin vanishing theorem,
affine scheme, Galois cohomology, cohomological dimension of fields, Tsen's
theorem, quasi-algebraically closed field, formally real field, Artin--Schreier
theory, fppf topology, finite group scheme, group of multiplicative type,
isotriviality, Tate cohomology, transfer, Sylow subgroup, inertia group,
decomposition group, Čech cohomology, Leray spectral sequence, injective sheaf,
flasque sheaf, Mayer--Vietoris, constructible sheaf, dévissage, Gysin sequence,
weak Lefschetz theorem, hyperplane section, Tate twist, local acyclicity, smooth
morphism, specialisation, Abhyankar's lemma, purity, simplicial action,
asphericity}
\end{resume}

\pagerange{1}{100}
\section*{Le fil du dossier, et les conventions}

\subsection*{Les stations}

Le dossier n'est pas ordonné, mais il est orienté. Voici ses stations, dans
l'ordre où le lecteur les rencontrera, et ce que chacune apporte à la question
centrale.

\begin{itemize}
\item \textbf{La dimension cohomologique d'un corps} — pages 2 à 8 (en
caractéristique $p$, pour la topologie fppf), 19 à 22 (les corps réels, où
elle est infinie), 46 à 51 (les corps de fonctions, où elle vaut la
dimension). C'est le socle : tout le reste s'y ramène par fibration.
\item \textbf{Les outils, et un contre-exemple} — pages 24 à 30 (quand les
groupes de Čech calculent la vraie cohomologie), 41 à 44 (comment une trace
inversible annule une cohomologie), 54 à 61 (pourquoi un faisceau injectif ne
reste pas acyclique sur un fermé : la réponse est \emph{non}, et c'est une
obstruction à la méthode évidente), 34 à 38 (asphéricité et torsion), 10 à 16
(un module dont un quotient cesse d'être de type fini).
\item \textbf{La dimension cohomologique des schémas affines} — pages 64 à 70
(la fonction $\varphi$, qui remplace un entier par une fonction sur les points),
73 à 80 (les deux théorèmes de finitude et leur récurrence croisée, dont sort
l'annulation d'Artin à la page 76), 81 à 84 (les conjectures qui restent).
\item \textbf{Deux conséquences géométriques} — pages 89 et 90 (Lefschetz
faible, qui \emph{consomme} l'annulation affine), 92 à 98 (la spécialisation et
l'asphéricité locale des morphismes lisses).
\item \textbf{Les feuillets de tapuscrit corrigés} — pages 11 à 17, 25 à 27,
71, 83 à 87. Ils ne sont pas de ce dossier-ci : ce sont des épreuves d'autres
exposés, sur les groupes de type multiplicatif et sur la complétion formelle,
qu'il a corrigées à la plume et dont il a utilisé les dos pour écrire le reste.
On les donne parce qu'ils portent sa main, et à part parce qu'ils ne sont pas
sur le même sujet.
\end{itemize}

Deux fils secondaires traversent le dossier sans s'y raccorder : un feuillet au
crayon sur les normes et l'imperfection (page 60) et une note sur une suite
spectrale (page 100). Ils sont donnés pour eux-mêmes à la fin.

\subsection*{Les conventions}

\emph{Schémas.} Le manuscrit dit « préschéma » là où l'on dit aujourd'hui
« schéma » ; il est antérieur au changement de vocabulaire d'EGA, où
« schéma » désignait encore le cas séparé. On écrit « schéma ».

\emph{Espaces affines.} Le manuscrit écrit $E^d_k$ pour l'espace affine de
dimension $d$ sur $k$, et $Y[t_1, \dots, t_r]$ pour l'espace affine relatif de
dimension $r$ sur $Y$. On écrit $\mathbf{A}^d_k$ et $\mathbf{A}^r_Y$.

\emph{Le nombre $\operatorname{cd}$.} Pour un schéma $X$ et un nombre premier
$\ell$, $\operatorname{cd}_{\ell}(X)$ est le plus petit $n$ tel que
$H^i(X,F) = 0$ pour $i > n$ et tout faisceau étale de $\ell$-torsion $F$ ; pour
un corps $k$, $\operatorname{cd}_{\ell}(k)$ est celui de son groupe de Galois
absolu. Sauf mention contraire $\ell$ est inversible sur le schéma considéré :
c'est l'hypothèse sans laquelle rien de ce qui suit ne vaut.\footnote{Le dossier
l'écrit tantôt « $\ell$ premier à la car.\ résiduelle », tantôt
« $\ell \cdot 1_X$ inversible », tantôt pas du tout. Aux endroits où elle est
omise mais nécessaire, la présente lecture la rétablit sans le signaler à
chaque fois.} Les pages 2 à 8 font exception et travaillent au contraire
\emph{en} la caractéristique, pour la topologie fppf ; on y revient là.

\emph{Cohomologie de Tate.} $\widehat{H}^q(G,M)$ désigne, pour $G$ fini, les
groupes de Tate, définis pour tout $q \in \mathbf{Z}$ ; $\widehat{H}(G,M) = 0$
sans exposant signifie : pour tout $q$.

\subsection*{Ce que le dossier annonce et n'établit pas}

Un chantier laisse des chantiers. Ceux-ci sont nommés ici une fois pour toutes,
pour qu'on ne les cherche pas dans le corps de la lecture :

\begin{itemize}
\item le Lemme 3 de la page 8, sur
$\operatorname{cd}^{\mathrm{add}}_p(k)$, s'arrête sur « se dévisse en groupes
$\alpha_p$ » et n'est pas démontré ;
\item à la page 42, les conditions b$'$), b($q$) et b$'$($q$) sont nommées et
leurs énoncés jamais écrits ;
\item la page 48 s'interrompt sur « Supposons que $k$ soit quelconque » ;
\item la démonstration du Th.\ 3 de la page 38 s'arrête à sa première ligne ;
\item le feuillet de tapuscrit de la page 25 casse au milieu du mot « à
l'aide » ;
\item la récurrence sur $r$ ouverte à la page 80 n'est pas conduite ;
\item la conjecture 1 bis) de la page 84 est écrite trois fois, chaque version
raturant la précédente, et n'est pas arrêtée ;
\item la Question marginale de la page 69 et celle de la page 22 (« cette
restriction est-elle essentielle ??? ») restent ouvertes dans le dossier.
\end{itemize}

\pagerange{2}{8}
\section*{La dimension cohomologique d'un corps}

\subsection*{En la caractéristique, pour la topologie fppf (pages 2 à 8)}

Les premières pages posent la question dans le cas où elle est le plus rétive :
$k$ un corps de caractéristique $p > 0$, et l'on veut borner la cohomologie à
coefficients de $p$-torsion. La cohomologie étale y est aveugle — un
$p$-groupe infinitésimal n'a pas de points étales — et il faut travailler dans
la topologie fidèlement plate de présentation finie.

On note $\operatorname{cd}^{!}_p(k)$ la dimension cohomologique calculée avec
les $k$-schémas en groupes commutatifs \emph{finis} d'ordre une puissance de
$p$. Tout tel groupe $G$ se dévisse en trois étages
\[
G' \subset G'' \subset G, \qquad G'' = G^{\circ},
\]
où $G/G''$ est étale, $G'$ est infinitésimal de type multiplicatif, et
$G''/G'$ infinitésimal de type additif. D'où :

\textbf{Lemme 1.} $\operatorname{cd}^{!}_p(k) \leqslant n$ si et seulement si
$H^i(k,G) = 0$ pour $i > n$ pour les $G$ de chacun des trois types : étale,
infinitésimal de type multiplicatif, infinitésimal de type
additif.\footnote{C'est le dévissage qui structure toute la suite. La page
l'écrit avec les abréviations « t.m.\ » et « add », et les trois dimensions
partielles $\operatorname{cd}^{\text{ét}}_p$,
$\operatorname{cd}^{\mathrm{t.m.}}_p$, $\operatorname{cd}^{\mathrm{add}}_p$.}

Le cas étale se règle d'un mot : $\operatorname{cd}^{\text{ét}}_p(k)$ vaut $0$
ou $1$ selon que $k$ n'a pas, ou a, une extension séparable finie de degré
divisible par $p$.

Le cas multiplicatif occupe les pages 4 à 8. Un groupe infinitésimal de type
multiplicatif est le dual d'un groupe fini étale, lequel est quotient d'un
module galoisien libre de type fini sur $\mathbf{Z}$ ; dualement, $G$ se plonge
dans un tore, et l'on a une suite exacte
\[
0 \longrightarrow G \longrightarrow T \longrightarrow T' \longrightarrow 0
\]
avec $T$, $T'$ des tores. D'où :

\textbf{Lemme 2.} $\operatorname{cd}^{\mathrm{t.m.}}_p(k) \leqslant n$ si et
seulement si, pour tout tore $T$ sur $k$, la composante $p$-primaire
$H^i(k,T)(p)$ est nulle pour $i > n$ et $p$-divisible pour $i = n$.
De façon équivalente, pour toute extension séparable finie $k'$ de $k$ :
$H^i(k', \mathbf{G}_m) = 0$ pour $i > n$, et $H^n(k', \mathbf{G}_m)$ est
$p$-divisible.

La nécessité se lit sur la suite de Kummer
$0 \to {}_pT \to T \xrightarrow{p} T \to 0$. Pour la suffisance, on descend à
une extension $k'$ de degré premier à $p$ : par le lemme de
Shapiro,\footnote{La page écrit $H^i(k, \prod_{k'/k} G') \simeq H^i(k,G)$ ;
l'énoncé de Shapiro donne $H^i(k, \prod_{k'/k} G') \simeq H^i(k', G')$, et
c'est ce dont l'argument a besoin. La lecture écrit la forme correcte.}
\[
H^i\Bigl(k, \prod\nolimits_{k'/k} G'\Bigr) \simeq H^i(k', G'),
\]
l'homomorphisme $\prod_{k'/k} G_{k'} \to G$ est un épimorphisme, et le
transfert permet de conclure sur $G$ dès qu'on conclut sur $G_{k'}$. On est
ainsi ramené au cas où $G$ est scindé par une $p$-extension, donc se dévisse en
copies de $\mu_p$, et tout revient à $H^{n+1}(\mu_p) = 0$, qui se lit sur
$0 \to \mu_p \to \mathbf{G}_m \to \mathbf{G}_m \to 0$.

Reste le cas additif, dont la page annonce la réponse sans la démontrer :

\textbf{Lemme 3.} $\operatorname{cd}^{\mathrm{add}}_p(k)$ vaut $0$ si $k$ est
parfait, $1$ sinon.\footnote{La démonstration s'arrête sur « tout groupe
infinitésimal additif se dévisse en groupes $\alpha_p$ » — la première ligne.
L'énoncé est celui de la page ; la présente lecture ne le certifie pas.}

\pagerange{19}{22}
\subsection*{Les corps réels, où la dimension est infinie (pages 19 à 22)}

Deux pages, écrites au net sous le titre « Théorèmes de M.\ Artin », enregistrent
la pathologie qui gouverne tout le reste du dossier : il existe des corps dont
la dimension cohomologique en $2$ est infinie, et l'on sait exactement
lesquels.

\textbf{Théorème 1.} Soit $K$ une extension de type fini de $\mathbf{Q}$. Les
conditions suivantes sont équivalentes :
\begin{enumerate}
\item[(i)] $-1$ est une somme de carrés dans $K$ ;
\item[(i bis)] $K$ n'admet aucune structure de corps ordonné ;
\item[(i ter)] $K$ n'est isomorphe à aucun sous-corps de $\mathbf{R}$ ;
\item[(ii)] il existe un modèle $X$ de $K$ sans point réel lisse ;
\item[(ii bis)] $\dim X(\mathbf{R}) < \deg\operatorname{tr} K/\mathbf{Q}$,
la dimension étant la dimension topologique ;
\item[(ii ter)] $X(\mathbf{R})$ est rare dans $X_{\mathbf{C}}$ ;
\item[(ii quater)] si $X$ est lisse sur $\mathbf{Q}$, $X(\mathbf{R}) =
\emptyset$ ;
\item[(iii)] $\operatorname{cd}_2(K) < +\infty$.
\end{enumerate}

L'équivalence des quatre premières est le théorème d'Artin--Schreier : un corps
est ordonnable si et seulement si $-1$ n'y est pas somme de carrés. Ce qui est
neuf est le raccord avec (iii) : un corps ordonnable a un élément d'ordre $2$
dans son groupe de Galois absolu — la conjugaison complexe d'une clôture
réelle — et un groupe profini ayant de la torsion a une dimension cohomologique
infinie.\footnote{Le dossier écrit « Théorèmes de M.\ Artin », c'est-à-dire
Michael Artin, son collaborateur de SGA 4 ; le cœur purement corporel de
l'énoncé est cependant classique et remonte à Artin et Schreier, et
l'observation sur la torsion et la dimension cohomologique est de Serre. La
lecture ne tranche pas la question d'attribution et se borne à nommer les
résultats.}

\textbf{Corollaire 1.} Soit $X$ un schéma de type fini sur $\mathbf{Z}$. Sont
équivalents : $X(\mathbf{R}) = \emptyset$ ; $\operatorname{cd}_2(X) < +\infty$ ;
$\operatorname{cd}_2(X_{\mathbf{Q}}) < +\infty$. Et alors
\[
\operatorname{cd}(X) < 2\dim X + 1, \qquad
\operatorname{cd}(X_{\mathbf{Q}}) \leqslant 2\dim X_{\mathbf{Q}} + 2, \qquad
\operatorname{cd}(X_{\mathbf{R}}) \leqslant 2\dim X.
\]

\textbf{Corollaire 3.} Pour $X$ de type fini sur $\mathbf{R}$, la finitude de
$\operatorname{cd}_2(X)$ équivaut encore à $X(\mathbf{R}) = \emptyset$, et l'on
a alors, pour tout faisceau étale $F$,
\[
H^i(X,F) \simeq H^i\bigl(X(\mathbf{C})/(\mathbf{Z}/2\mathbf{Z}),\, F\bigr)
\]
— la cohomologie ordinaire du quotient de l'espace des points complexes par la
conjugaison — et $\operatorname{cd}_2(X) \leqslant 2\dim X$.

Vient ensuite l'énoncé qui servira au dossier tout entier, et qui dit combien
de cohomologie on gagne en montant en degré de transcendance :

\textbf{Théorème 2.} Soient $k$ un corps et $r$ un nombre premier. Sont
équivalents :
\begin{enumerate}
\item[(i)] pour toute extension de type fini $k'$ de $k$,
$\operatorname{cd}_r(k') = \operatorname{cd}_r(k) + \deg\operatorname{tr}
k'/k$ ;
\item[(i bis)] la même égalité pour la seule extension
$k' = k(\sqrt{-1})$ ;
\item[(iii)] on n'est pas dans le cas d'exception : $k$ ordonnable (donc de
caractéristique $0$), $r = 2$, et $\operatorname{cd}_2(k(\sqrt{-1})) <
+\infty$.
\end{enumerate}
La marge énumère les cas où la condition est acquise sans avoir à la vérifier :
$\operatorname{car} k > 0$ ; $r \neq 2$ ; $-1$ somme de carrés dans $k$ ;
$\operatorname{cd}_2 k(\sqrt{-1}) < +\infty$.\footnote{La condition (ii) de la
page — l'existence d'un sous-corps $k_0$ de type fini sur le corps premier avec
$\operatorname{cd}_r(k_0)$ fini — est écrite avec un $k'$ là où le sens demande
$k$, et n'est pas reprise ici.}

\textbf{Corollaire.} Soit $A$ un anneau local noethérien intègre, de corps des
fractions $K$ et de corps résiduel $k$. Alors, pour tout premier $r$,
\[
\operatorname{cd}_r(K) \geqslant \operatorname{cd}_r(k) + \dim A,
\]
à deux exceptions possibles près : $r$ égale la caractéristique de $K$ ; ou
bien $r = 2$ et il existe $x \in \operatorname{Spec}(A)$ dont le corps résiduel
$k(x)$ est ordonnable avec $\operatorname{cd}_2 k(x)(\sqrt{-1})$ fini. La marge
donne l'exemple qui montre que le second cas se produit :
$A = \mathbf{R}[X,Y]/(X^2+Y^2)$ localisé. Si $A$ est caténaire et
universellement japonais, l'exception se précise : c'est $k$ lui-même qui doit
être réel. Si $A$ est régulier, elle disparaît et l'inégalité vaut sans
réserve.

Et, dans la marge, en travers : \emph{cette restriction est-elle essentielle
???} Le dossier ne répond pas.

\pagerange{46}{51}
\subsection*{Les corps de fonctions, où la dimension vaut la dimension (pages 46 à 51)}

Quinze pages plus loin, le dossier reprend la même question par l'autre bout, du
côté où la réponse est bonne. La Proposition de la page 46 caractérise
l'annulation de la cohomologie galoisienne au-delà d'un rang $n$ par quatre
conditions de plus en plus maniables : elle vaut pour tous les groupes
algébriques commutatifs si et seulement si elle vaut pour les groupes finis, si
et seulement si elle vaut pour les groupes des éléments inversibles des
algèbres commutatives séparables.\footnote{L'énoncé de la page 46 est un
brouillon rapide : les conditions b$'$) et c) sont écrites par-dessus des
versions biffées et la prose qui les relie est en grande partie perdue. La
lecture donne la charpente — quatre conditions, et le circuit d'implications
a) $\Rightarrow$ b), a) $\Rightarrow$ c), c) $\Rightarrow$ b$'$),
b$'$) $\Rightarrow$ b) — et ne reconstitue pas les énoncés que le feuillet ne
porte plus.}

Les dévissages qui la démontrent sont instructifs et tiennent en trois gestes.
Pour passer des groupes finis à tous les groupes algébriques commutatifs : on
sépare la partie unipotente, on utilise que la décomposition en parties
semi-simple et unipotente est définie sur $k$, et l'on se ramène à trois cas —
unipotent, à éléments semi-simples, commutatif complet. Dans les deux derniers,
on passe au sous-groupe $\mathcal{U}$ des éléments d'ordre fini : le quotient
$G/\mathcal{U}$ est divisible, donc sans cohomologie, d'où
$H^i(\mathfrak{g}, G) \simeq H^i(\mathfrak{g}, \mathcal{U})$ pour $i \geqslant
1$, et $\mathcal{U}$ est réunion filtrante de ses sous-groupes finis
${}_n\mathcal{U}$, où l'hypothèse s'applique.\footnote{La page passe sans
prévenir de $H^i(k,-)$ à $H^i(\mathfrak{g},-)$, $\mathfrak{g}$ désignant le
groupe de Galois absolu ; on a gardé la distinction. Une note marginale de sa
main, à cet endroit, dit : « ici c'est faux », sans qu'on puisse établir à
quelle ligne elle se rapporte.} Pour le cas unipotent, on distingue la
caractéristique nulle — où $G \simeq \mathbf{G}_a^n$ déjà sur $k$, car
$H^1(k, \mathrm{GL}(n)) = 0$ — et la caractéristique $p$, où la filtration
$G \supset pG \supset \cdots$ a tous ses crans définis sur $k$ et ramène à
$pG = 0$.

Puis les deux lemmes qui donnent la réponse pour les corps de fonctions.

\textbf{Lemme 1 (Tsen).} Soient $k$ algébriquement clos et $K$ une extension de
type fini de $k$, de degré de transcendance $1$. Alors $K$ est
quasi-algébriquement clos : toute forme homogène de degré $d$ en plus de $d$
variables y a un zéro non trivial.

\textbf{Corollaire 1.} $H^i(K, \mathbf{G}_m) = 0$ pour $i \geqslant 1$. En
particulier le groupe de Brauer de $K$ est nul.\footnote{La page écrit
$\widehat{H}^i(K, \mathbf{G}_m) = 0$ « pour tout $i$ », en cohomologie de Tate.
Pour un groupe profini les groupes de Tate en degré $\geqslant 1$ coïncident
avec les groupes ordinaires, et c'est de cette forme que la suite se sert ; le
degré $0$ n'est pas nul en général. La lecture écrit donc l'énoncé en degrés
$\geqslant 1$.}

\textbf{Corollaire 2.} $H^i(K, \Gamma) = 0$ pour $i > 1$ et tout module
galoisien fini $\Gamma$ ; autrement dit
$\operatorname{cd}(K) \leqslant 1$.\footnote{La page énonce le corollaire pour
$\Gamma$ fini à action triviale. L'argument qu'elle donne — la surjectivité de
$x \mapsto x^n$ sur $\widetilde{K}^{*}$ pour $n$ premier à la caractéristique,
et la suite d'Artin--Schreier
$0 \to \mathbf{Z}/p \to \widetilde{K} \xrightarrow{\wp} \widetilde{K} \to 0$
pour la partie $p$ — donne la majoration $\operatorname{cd}(K) \leqslant 1$,
d'où l'annulation pour \emph{tout} module galoisien fini. La lecture énonce
cette forme, qui est celle dont le Lemme 2 a besoin.}

\textbf{Lemme 2.} Soit $k$ un corps tel que $H^i(k', \Gamma) = 0$ pour $i > n$
et toute extension de type fini $k'$ de $k$. Alors $H^i(K, \Gamma) = 0$ pour
$i > n+1$ dès que $K$ est de degré de transcendance $\leqslant 1$ sur $k$.

La démonstration est le mécanisme même de tout le dossier, dans sa forme la plus
simple. On intercale entre $k$ et $K$ le composé avec une clôture algébrique,
d'où une suite exacte de groupes de Galois
\[
e \longrightarrow G' \longrightarrow G \longrightarrow G'' \longrightarrow e
\]
et sa suite spectrale
$H^i\bigl(G'', H^j(G', \Gamma)\bigr) \Rightarrow H^{*}(G, \Gamma)$. Le Lemme 1
tue les $j > 1$ ; l'hypothèse tue les $i > n$ ; il ne reste rien au-dessus de
$n+1$.

\textbf{Corollaire.} Si $K$ est un corps de fonctions algébriques de dimension
$n$ sur un corps algébriquement clos, $\operatorname{cd}(K) \leqslant n$.

C'est le socle annoncé. Il dit que, du côté géométrique, un corps se comporte
comme sa dimension le laisse espérer ; toute la difficulté restante est de
passer d'un corps à un schéma.

\pagerange{24}{30}
\section*{Les outils, et un contre-exemple}

\subsection*{Le lemme de Cartan : quand Čech suffit (pages 24 à 30)}

Pour calculer une cohomologie il faut une résolution injective, qu'on n'a
jamais. Ce qu'on a, ce sont des recouvrements et les groupes de Čech qu'ils
produisent, calculables à la main. La question de savoir quand les seconds
donnent la première est donc préalable à tout, et le dossier la règle sous un
nom que Grothendieck lui donne lui-même — le lemme de Cartan.

Soient $T$ un site, $F$ un faisceau abélien sur $T$, $n$ un entier, et
$\mathcal{M}$ une classe d'objets de $T$ satisfaisant deux conditions :
$\mathcal{M}$ est stable par produits fibrés, et pour tout $S \in \mathcal{M}$
les recouvrements de $S$ par des objets de $\mathcal{M}$ sont cofinaux parmi
tous les recouvrements de $S$.

\textbf{Lemme de Cartan.} Sous ces conditions, les énoncés suivants sont
équivalents, chacun pour $0 < i \leqslant n$ :
\begin{enumerate}
\item[(i)] $H^i(S,F) = 0$ pour tout $S \in \mathcal{M}$ ;
\item[(ii)] $\check{H}^i(\mathcal{U},F) = 0$ pour tout $S \in \mathcal{M}$ et
tout recouvrement $\mathcal{U}$ de $S$ par des objets de $\mathcal{M}$ ;
\item[(iii)] $\varinjlim_{\mathcal{U}} \check{H}^i(\mathcal{U},F) = 0$,
la limite portant sur ces mêmes recouvrements.
\end{enumerate}

Tout passe par une seule suite spectrale, celle qui compare Čech au foncteur
dérivé :
\[
E_2^{pq} = \check{H}^p\bigl(\mathcal{U}, \mathcal{H}^q(F)\bigr)
\Longrightarrow H^{p+q}(S,F),
\]
où $\mathcal{H}^q(F)$ est le préfaisceau $U \mapsto H^q(U,F)$. Sous (i) les
lignes $0 < q \leqslant n$ sont nulles, d'où $\check{H}^i(\mathcal{U},F)
\simeq H^i(S,F)$ pour $i \leqslant n$ et donc (ii) ; (ii) $\Rightarrow$ (iii)
est trivial ; et (iii) $\Rightarrow$ (i) se fait par récurrence sur $n$, en
observant qu'à l'étape $n$ l'hypothèse de récurrence annule les colonnes
$0 < q < n$, l'hypothèse (iii) la ligne $q = 0$ pour $0 < p \leqslant n$, et la
définition d'un faisceau la colonne $p = 0$ pour $q > 0$ : il ne reste aucun
terme sur la diagonale $p+q = n$.\footnote{La page accompagne l'argument d'un
croquis du plan $(p,q)$ où l'on voit la diagonale ne rencontrer aucune croix ;
c'est la figure qui porte la démonstration, et elle n'a pas d'autre contenu que
ce décompte.}

La page 28 pousse plus loin et demande non plus l'annulation mais
l'\emph{identification} : quand a-t-on
$\varinjlim_{\mathcal{U}} \check{H}^i(\mathcal{U},F) \simeq H^i(S,F)$ pour
$i \leqslant n$ ? On introduit pour cela le préfaisceau
\[
\overline{H}^q(F)(T) = \varinjlim_{\mathcal{V}}
\check{H}^q(\mathcal{V}, F),
\]
$\mathcal{V}$ parcourant les recouvrements de $T$ par des objets de
$\mathcal{M}$ — autrement dit la version « déjà passée à la limite » de
$\mathcal{H}^q(F)$. La condition suffisante est l'annulation de
$\varinjlim_{\mathcal{U}} \check{H}^p(\mathcal{U}, \mathcal{H}^q(F))$, et de
même avec $\overline{H}^q(F)$, pour $q > 0$ et $p+q \leqslant
n$.\footnote{L'intervalle d'indices est, sur le feuillet, réécrit par-dessus
lui-même : on lit « $p < q \leqslant n$ » là où une version biffée portait
« $0 < q \leqslant n$ », et une troisième ligne, biffée aussi, porte
« $p + q \leqslant n-1$ ». La lecture donne l'intervalle dont la démonstration
de la page 30 se sert effectivement — $q > 0$ et $p+q \leqslant n$ — et signale
que le feuillet ne le fixe pas.} Deux variantes au niveau des cochaînes, (iii
bis) et (iii ter), disent la même chose de façon vérifiable : toute cochaîne de
$C^p(\mathcal{U}, \mathcal{H}^q(F))$ s'annule sur un raffinement.

La page 30 démontre le circuit, de nouveau par récurrence sur $n$ et de nouveau
sur la même suite spectrale, et le referme par un diagramme :

\begin{tikzcd}[column sep=small, row sep=small]
\text{(ii)} \arrow[d, Rightarrow] \arrow[dr, Rightarrow] & & \text{(iii bis)} \arrow[ll, Rightarrow] \arrow[d, Rightarrow] \\
\text{(i)} & \text{(iii)} \arrow[l, Rightarrow] & \text{(iii ter)} \arrow[l, Rightarrow]
\end{tikzcd}

\pagerange{41}{44}
\subsection*{La trace, ou comment tuer une cohomologie (pages 41 à 44)}

Le deuxième outil est un principe d'une simplicité désarmante, et c'est lui qui
justifie toutes les hypothèses « $\ell$ premier à » du dossier.

Soient $B$ un anneau unitaire, $G$ un groupe fini d'automorphismes de $B$, et
$A = B^G$. On pose $\operatorname{Tr}_G b = \sum_{g \in G} gb$.

\textbf{Proposition.} Sont équivalents :
\begin{itemize}
\item[a)] $A = \operatorname{Tr}_G B$ ;
\item[a$'$)] $1 \in \operatorname{Tr}_G B$ ;
\item[b)] $\widehat{H}(G,M) = 0$ pour tout $(G,B)$-module $M$ ;
\item[b($q$))] la même annulation en un seul degré $q$.
\end{itemize}

Le ressort est une ligne. Si $\varphi_M(b)$ désigne l'homothétie de rapport $b$
dans $M$, on a $g\,\varphi_M(b)\,g^{-1} = \varphi_M(gb)$, donc
$\varphi_M(\operatorname{Tr}_G b) = \operatorname{Tr}_G \varphi_M(b)$ ; et une
homothétie dont le rapport est une trace agit par zéro sur la cohomologie de
Tate. Si donc $1$ est une trace, l'identité de $\widehat{H}(G,M)$ est nulle, et
le groupe l'est aussi. Le passage d'un degré au suivant se fait sur la suite
exacte
\[
0 \longrightarrow M \longrightarrow M \otimes_{\mathbf{Z}} \mathbf{Z}[G]
\longrightarrow N \longrightarrow 0, \qquad m \longmapsto \sum_g m \otimes g,
\]
dont le terme du milieu est induit, donc de cohomologie de Tate nulle : d'où
$\widehat{H}^{q+1}(G,M) \simeq \widehat{H}^{q}(G,N)$ et la récurrence dans les
deux sens.

La page 42 en donne la forme relative, pour $G' \subset G$ et $A' = B^{G'}$,
avec le transfert $t(G,G')$. Et elle énonce le théorème dont tout le reste
descend :

\textbf{Théorème.} Soient $G$ fini, $G' \subset G$, et $M$ un $G$-module tel
que l'identité de $M$ s'écrive $1 = \operatorname{Tr}_{G/G'}\varphi$ pour un
$\varphi \in \operatorname{Hom}_{G'}(M,M)$. Alors le transfert
$\widehat{H}(G',M) \to \widehat{H}(G,M)$ est surjectif, la restriction
$\widehat{H}(G,M) \to \widehat{H}(G',M)$ est injective, et son image est le
sous-groupe des éléments stables $\widehat{H}(G',M)^{G}$.

Suit, entre parenthèses et de sa main : « Ceci contient le Th.\ sur les groupes
de Sylow ». C'est exact — en prenant pour $G'$ un $p$-Sylow, l'indice $[G:G']$
est premier à $p$ et l'hypothèse est acquise sur la composante $p$-primaire.
La marge indique l'autre application visée : prendre $G' = G_i$, le groupe
d'inertie.

C'est ce que fait la page 43. Soit $G$ un groupe fini d'automorphismes d'un
anneau local $\mathcal{O}$, transitif sur les idéaux maximaux ; soient $G_d$ le
groupe de décomposition de $\mathfrak{m}_0$ et $G_i \subset G_d$ le groupe
d'inertie, noyau de l'action sur le corps résiduel. S'il existe
$x \in \mathcal{O}^{G_i}$ avec $\operatorname{Tr}_{G/G_i} x = 1$, alors pour
tout module $M$
\[
\widehat{H}(G,M) = \widehat{H}(G_i,M)^{G/G_i},
\]
et $\widehat{H}(G,M) = 0$ pour tout $M$ si et seulement si l'ordre de $G_i$
n'est pas divisible par la caractéristique résiduelle $p$. C'est l'énoncé
classique : \emph{une extension modérément ramifiée est cohomologiquement
triviale}. En particulier, si $G_i$ est trivial — extension non ramifiée —
$H^1(G, \mathcal{O}^{*}) = \{e\}$.

La page se ferme sur un carré de comparaison entre $\widehat{H}(G,M)$,
$\widehat{H}(G_i,M)^{G/G_i}$, $\widehat{H}(G_d,M)^{G/G_d}$ et
$\widehat{H}(G_d, M_{\mathfrak{m}_0})$, dont deux flèches sont annotées en
marge : « injectifs, à images denses ». Et, à côté du corollaire :
\emph{Comment voir ça ?}

\pagerange{54}{61}
\subsection*{Un faisceau injectif ne reste pas acyclique sur un fermé (pages 54 à 61)}

Voici la page où le dossier essuie un refus, et c'est l'une des plus utiles.

La méthode évidente pour calculer la cohomologie d'un fermé $Y \subset X$ est
de prendre une résolution injective sur $X$ et de la restreindre. Elle marche
si la restriction d'un injectif est acyclique. La page 54 pose la question :
$X$ un espace topologique, $Y$ un fermé, $F$ un faisceau injectif sur $X$, a-t-on
$H^i(Y, F|Y) = 0$ pour $i > 0$ ?

\emph{Réponse négative.}

La raison de fond est formelle et vaut la peine d'être dite, car la page la
laisse implicite. Pour une immersion fermée $i : Y \to X$, l'image directe
$i_{*}$ est exacte et admet $i^{-1}$ pour adjoint à gauche ; c'est pourquoi
$i_{*}$ préserve les injectifs. Mais $i^{-1}$ n'a pas d'adjoint à gauche exact
— l'extension par zéro est adjointe à gauche de $i^{-1}$ pour une immersion
\emph{ouverte}, pas fermée — et rien ne fait donc de $i^{-1}F$ un injectif, ni
même un acyclique.

Le dossier ne raisonne pas ainsi : il construit. Et il commence par montrer que
la chose casse déjà au degré $0$.

\textbf{Exemple 1.} Soit $X$ un espace irréductible et $Y = \{a,b\}$ un fermé
formé de deux points fermés distincts. Alors $Y$ est discret, et une section de
$F|Y$ vaut $0$ en $a$ et $1$ en $b$ sans provenir d'aucune section sur un
voisinage de $Y$ : car un tel voisinage contiendrait un ouvert autour de $a$ et
un ouvert autour de $b$, et dans un espace irréductible deux ouverts non vides
se rencontrent. L'homomorphisme
\[
\varinjlim_{U \supset Y} H^0(U,F) \longrightarrow H^0(Y, F|Y)
\]
n'est donc pas surjectif, quelque injectif que soit $F$.\footnote{C'est le point
précis où l'intuition venue des espaces séparés trompe : pour $Y$ compact dans
un espace raisonnable, les sections de $i^{-1}F$ sur $Y$ \emph{sont} la limite
inductive des sections sur les voisinages. La topologie de Zariski n'est pas de
ce genre, et l'irréductibilité est exactement ce qui casse l'énoncé.}

\textbf{Exemples 2 et 3} varient le motif — un cercle coupé par une droite en
deux points, puis deux droites se coupant en un point, avec le faisceau des
fonctions à valeurs dans $\mathbf{Q}$ — et montrent que passer à un $Y$
irréductible ne suffit pas à sauver la situation.

\textbf{Exemple 4} donne enfin le contre-exemple à la question elle-même. Soit
$Y = Y_1 \cup Y_2$, réunion de deux courbes irréductibles se coupant en deux
points $a$ et $b$, dans un $X$ irréductible ; $F$ le faisceau des fonctions de
$X$ dans $\mathbf{Q}$. La suite de Mayer--Vietoris donne
\[
H^0(Y_1, F_1) \times H^0(Y_2, F_2) \longrightarrow H^0(Y_{12}, F_{12})
\longrightarrow H^1(Y, F|Y) \longrightarrow 0
\]
avec $H^0(Y_{12}, F_{12}) = F_a \times F_b$. L'élément $(0,1)$ ne provient pas
du terme de gauche : s'il en provenait, deux sections $\varphi_1, \varphi_2$
coïncideraient au voisinage de $a$ et différeraient de $1$ au voisinage de $b$,
et les deux voisinages se rencontrent puisque $X$ est irréductible — d'où
$0 = 1$. Donc $H^1(Y, F|Y) \neq 0$ pour un $F$ injectif, ce qu'il fallait
produire. La page 58 relève l'exemple en dimension $2$, avec un $Y$ irréductible,
en coupant $X$ par un plan qui recoupe $Y$ en deux composantes.

La Proposition des pages 59 et 61 tire la morale positive : pour un faisceau
abélien $F$ sur un espace noethérien, l'annulation de $H^i(U,F)$ pour tout
ouvert $U$ et tout $i > 0$, l'annulation des groupes de Čech de tout
recouvrement, et la flasquité sont équivalentes.\footnote{Les conditions de la
page 59 sont numérotées (i), (i bis), (ii), (ii bis), (iii) et leurs énoncés
n'y sont écrits qu'à demi ; la présente lecture ne restitue que celles que la
page porte en clair et que la démonstration de la page 61 utilise.} La dernière
implication est celle que la page 61 achève : on plonge $F$ dans un flasque
$C^0$, l'hypothèse fait de la suite une suite exacte de flasques, donc $C^1$
est flasque comme $C^0$, donc la résolution calcule la cohomologie, et les
$H^i(U,F)$ s'annulent au-delà du premier cran par décalage.

\pagerange{34}{38}
\subsection*{Asphéricité et torsion (pages 34 à 38)}

Cinq feuillets à l'encre bleue, d'une main qui n'est pas la sienne et sous
l'en-tête « GRÖTHENDIECK. », enregistrent un résultat d'un genre différent : il
n'y est question ni de schémas ni de cohomologie étale, mais de groupes opérant
sur des ensembles simpliciaux. Il est là parce qu'il fournit le mécanisme par
lequel un quotient par un groupe fini ne fait apparaître que de la torsion
contrôlée — le même mécanisme, au fond, que la trace de la page 41.

\textbf{Théorème 1.} Soient $E$ un ensemble simplicial connexe, $G$ un groupe
discret opérant sur $E$ par transformations simpliciales, $B = E/G$ le quotient
et $p : E \to B$ la projection. Alors l'image de
$p_{*} : \pi_1(E) \to \pi_1(B)$ est distinguée, et
\[
\pi_1(B)/\operatorname{im} p_{*} \simeq G/S,
\]
où $S$ est le sous-groupe distingué de $G$ engendré par les stabilisateurs des
sommets ; $E/S$ est le revêtement de $B$ correspondant à
$\operatorname{im} p_{*}$.

La démonstration construit directement l'isomorphisme : un $g \in G$ étant
donné, on choisit un chemin de $x_0$ à $g x_0$, sa projection est un lacet de
$B$, et sa classe modulo $\operatorname{im}\pi_1(E)$ ne dépend pas du chemin.
L'application $\chi$ ainsi obtenue est surjective (les chemins se relèvent,
Lemme 1), multiplicative, et son noyau est exactement $S$.

\textbf{Théorème 2.} Si $G$ est fini d'ordre $n$, l'image de
$H_{*}(E) \to H_{*}(B)$ contient tous les éléments divisibles par $n$.

La preuve tient en une formule : pour un cycle $\zeta = \sum a_i \sigma_i$ de
$B$, on relève chaque simplexe et l'on moyenne sur le groupe ; le cycle
$\theta = \sum_{g,i} g\,a_i \tau_i$ de $E$ se projette sur $n\zeta$. C'est
encore la trace, sous un autre habit.

\textbf{Théorème 3.} Si $E$ est contractile et $S$ fini d'ordre $n$, alors pour
$i > 1$ tout élément de $\pi_i(B)$ est d'ordre divisant une puissance de
$n$.\footnote{La démonstration s'arrête à sa deuxième ligne : $E/S$ est le
revêtement universel de $B$ par le Théorème 1, donc $\pi_i(B) = \pi_i(E/S)$, et
l'on veut conclure par le Théorème 2 sur les $H_i(E/S)$. Le feuillet s'arrête
là.}

\pagerange{10}{16}
\subsection*{Un module dont un quotient cesse d'être de type fini (pages 10 à 16)}

Quatre feuillets, écrits au dos du tapuscrit, construisent un contre-exemple de
finitude. Il n'appartient pas à la ligne principale mais il en éclaire les
précautions : les théorèmes d'annulation du dossier portent tous sur des
faisceaux constructibles ou noethériens, et voici ce qui arrive sans cette
hypothèse.

Soient $A$ un anneau, $f \in A$ un élément $A$-régulier, $g \in A$ un élément
$B$-régulier où $B = A/fA$, avec $B/gB \neq 0$ — par exemple $A$ local régulier
de dimension $2$ et $(f,g)$ un système de paramètres. On construit un
$A$-module $M$ tel que
\begin{enumerate}
\item[(i)] $M = \bigcup_{n} {}_{f^n}M$ avec
${}_{f^{n+1}}M/{}_{f^n}M \simeq B$ — donc $M$ est de $f$-torsion, à support
dans $V(f)$, et ses crans successifs sont aussi petits qu'on veut ;
\item[(ii)] $M/fM \simeq (B/gB)^{(\mathbf{N})}$ — donc $M/fM$ n'est pas de type
fini.
\end{enumerate}

La construction est une récurrence sur des extensions. On pose $E_{n+1}$ égale
à l'image directe (le \emph{pushout}) de la suite
$0 \to A \xrightarrow{f} A \to A/fA \to 0$ le long de l'homomorphisme
$A \to E_n$ défini par un élément $\xi_n \in E_n$ :

\begin{tikzcd}[column sep=small, row sep=small]
0 \arrow[r] & A \arrow[r, "f"] \arrow[d, "\xi_n"'] & A \arrow[r] \arrow[d] & A/fA \arrow[r] \arrow[d, no head, "\|" description] & 0 \\
0 \arrow[r] & E_n \arrow[r] & E_{n+1} \arrow[r] & A/fA \arrow[r] & 0
\end{tikzcd}

L'extension ne dépend que de $\xi_n$ modulo $fE_n$, et tout le calcul consiste à
choisir $\xi_n = (g,0)$ — c'est le $g$ injecté à chaque cran qui fait grossir
le quotient — puis à vérifier les trois conditions de récurrence :
$f^n E_n = 0$ ; ${}_{f^{k+1}}E_n/{}_{f^k}E_n \simeq B$ ; et
$E_n/fE_n \simeq B \times (B/gB)^{n-1}$. La vérification de la troisième est le
calcul qui porte tout :
\[
E_{n+1}/fE_{n+1} \simeq
\bigl([B \times (B/gB)^{n-1}] \times B\bigr) \big/ B(g,0,0)
\simeq B \times (B/gB)^{n}.
\]
On pose $M = \varinjlim E_n$, et le quotient $M/fM$ est la somme directe
infinie annoncée.\footnote{La prose des pages 12, 14 et 16 est écrite très vite
et en grande partie perdue ; les formules, elles, passent. La lecture donne la
construction et le calcul, et ne reconstitue pas les liaisons.}

\pagerange{64}{70}
\section*{La dimension cohomologique des schémas affines}

\subsection*{Faisceaux élémentaires et la fonction $\varphi$ (pages 64 à 70)}

On arrive au cœur. Les pages 64 à 70, annoncées par une chemise portant de sa
main « cd des affines », introduisent le dispositif qui remplace un \emph{entier}
par une \emph{fonction sur les points} — et c'est ce raffinement qui rend la
récurrence possible.

Le point de départ est un dévissage. Appelons \emph{$\ell$-élémentaire} un
faisceau de la forme $i_{*}(G)_{U,X}$, où $x$ est un point de $X$,
$i : \operatorname{Spec} k(x) \to X$ l'immersion canonique, $G$ un faisceau fini
de $\ell$-torsion, et $U$ un ouvert de $Y = \overline{\{x\}}$ ; il est dit
\emph{strictement} élémentaire lorsque $G$ est étale sur $U$ et $U$ unibranche.

\textbf{Proposition.} Tout faisceau de $\ell$-torsion noethérien sur un schéma
noethérien $X$ admet une suite de composition finie dont les quotients sont
strictement élémentaires.

C'est le dévissage des faisceaux constructibles, et il ramène toute question
d'annulation à ces briques-là. D'où :

\textbf{Corollaire 1.} $\operatorname{cd}_{\ell}(X) \leqslant n$ si et seulement
si, pour tout fermé intègre $Y \subset X$, tout fermé $Z \subset Y$, tout point
$x$ et tout faisceau fini de $\ell$-torsion $G$ sur $\operatorname{Spec} k(x)$ :
\begin{itemize}
\item[a)] $H^n(Y, i_{*}G) \to H^n(Z, i_{*}G|Z)$ est surjectif,
\item[b)] $H^{n+1}(Y, i_{*}G) = 0$.
\end{itemize}
Une variante (ii bis) remplace l'immersion par un morphisme fini $Y \to X$ avec
$Y$ intègre unibranche ; et un Corollaire 1 bis dit la même chose pour la
cohomologie à support dans un fermé $T$.

Vient alors le raffinement. On remplace l'entier $n$ par une fonction
$\varphi : X \to \mathbf{N}$, et l'on pose pour toute partie fermée $Y$
\[
\varphi(Y) = \operatorname*{Sup}_{x \in Y} \varphi(x)
= \operatorname*{Max}_{x \text{ maximal dans } Y} \varphi(x),
\]
la seconde égalité étant ce qui rend la fonction maniable : elle ne se lit qu'aux
points maximaux. Le Corollaire 2 transporte le critère précédent mot pour mot,
avec $\varphi(Y)$ à la place de $n$.

Le théorème que tout cela prépare est celui-ci :

\textbf{Proposition (page 66).} Soit $\varphi$ une fonction sur $X$ telle que
\begin{itemize}
\item[a)] $\varphi(x) \geqslant \operatorname{cd}_{\ell}(k(x))$ pour tout
$x \in X$ ;
\item[b)] pour toute spécialisation $y \in \overline{\{x\}}$ avec $y \neq x$,
\[
\varphi(y) < \varphi(x) - \operatorname{cd}_{\ell}
\bigl(\mathcal{O}_{X,\bar{y}} \otimes k(x)\bigr),
\]
$\mathcal{O}_{X,\bar{y}}$ désignant l'hensélisé strict.
\end{itemize}
Alors, pour tout faisceau $F$ de $\ell$-torsion,
$H^i(X,F) = 0$ dès que $i > \varphi(F) =
\operatorname*{Sup}_{x \in \operatorname{Supp} F} \varphi(x)$.

Il faut lire ces deux conditions pour ce qu'elles sont. La première dit que
$\varphi$ majore ce que les corps résiduels apportent déjà — c'est le socle
établi plus haut. La seconde est une condition de \emph{décroissance stricte le
long des spécialisations}, et le terme qu'on retranche est exactement la
cohomologie de l'anneau local strict relatif $K_{\bar{y}} =
\mathcal{O}_{X,\bar{y}} \otimes k(x)$ : en descendant d'un point à une de ses
spécialisations, $\varphi$ doit perdre au moins ce que la fibre peut faire
remonter.

La démonstration est une récurrence noethérienne. On se ramène à
$F = g_{*}(G)$ avec $g : \overline{\{x\}} \to X$, on écrit la suite spectrale
\[
E_2^{pq} = H^p\bigl(X, R^q g_{*}(G)\bigr) \Longrightarrow
H^{p+q}\bigl(\overline{\{x\}}, G\bigr),
\]
et tout tient dans une majoration, encadrée sur la page :
\[
\varphi\bigl(R^q g_{*}(G)\bigr) \leqslant \varphi(x) - q.
\]
Elle s'obtient en calculant la fibre $R^q g_{*}(G)_{\bar{y}} =
H^q(K_{\bar{y}}, G)$, qui est nulle dès que $q > \operatorname{cd}_{\ell}
(K_{\bar{y}})$, et en reportant dans la condition b). L'hypothèse de récurrence
donne alors $E_2^{pq} = 0$ pour $p + q > \varphi(x)$, et il ne reste rien
au-dessus.

La page 68 note au passage la façon dont $\varphi$ se comporte le long d'un
morphisme de type fini $f : X \to Y$ : pour $x \in X$ et $y = f(x)$,
\[
\varphi(x) \leqslant \varphi(y) + 2\deg\operatorname{tr}\bigl(k(x)/k(y)\bigr).
\]
Le facteur $2$ est celui du Théorème 2 de la page 21 : chaque degré de
transcendance coûte deux crans de dimension cohomologique, pas un.

\textbf{Corollaire (page 69).} Supposons que partout
$\operatorname{cd}_{\ell} K_{x,\bar{y}} \leqslant \dim \mathcal{O}_{X,\bar{y}}$.
Alors la fonction
\[
\varphi(x) = \operatorname*{Sup}_{y \in \overline{\{x\}}}
\Bigl(\operatorname{cd}_{\ell}\bigl(k(y)\bigr)
+ 2\dim \mathcal{O}_{\overline{\{x\}},y}\Bigr)
\]
satisfait a) et b), on a pour toute partie fermée $Y$ la même formule avec
$\mathcal{O}_{Y,y}$, et $\operatorname{cd}_{\ell}(X) \leqslant \varphi(X)$.

La page en donne quatre domaines d'application : les schémas algébriques et plus
généralement les schémas de type fini ; les corps de nombres ; le spectre d'un
anneau strictement local ; les schémas affines de dimension $1$.

Et dans la marge, encadrée, la question qui reste :

\begin{quote}
Soit $A$ local noethérien hensélien intègre, de corps des fractions $K$ et de
corps résiduel $k$ avec $\operatorname{cd}_{\ell}(k) < +\infty$. A-t-on
$\operatorname{cd}_{\ell} K \geqslant \dim A + \operatorname{cd}_{\ell} k$ ?
\end{quote}

C'est la réciproque du Corollaire de la page 22, et le dossier ne la tranche
pas.

\pagerange{73}{80}
\subsection*{Les deux rungs de la récurrence, et le théorème d'Artin (pages 73 à 80)}

Huit pages écrites à toute vitesse, dont la prose de liaison est en grande partie
perdue et dont les énoncés, eux, passent. Elles portent la démonstration
proprement dite, et son architecture est une \emph{récurrence croisée} entre
deux familles d'énoncés.

Les deux familles sont indexées par un entier $n$, et le dossier les note
$1_n$ et $2_n$ (puis $1'_n$, $2'_n$ pour les variantes restreintes aux schémas
de type fini sur un corps) :

\textbf{Énoncé $1_n$.} Pour $A$ strictement local noethérien de dimension $n$ et
$F$ un faisceau de $\ell$-torsion « nul en codimension $< m$ » — c'est-à-dire
tel que $x \in \operatorname{Supp} F$ entraîne
$\dim \mathcal{O}_{X,x} \geqslant m$ — on a $H^i(X,F) = 0$ pour
$i > n+1-m$.\footnote{Tel que la page l'écrit, avec $X = \operatorname{Spec} A$,
l'énoncé est vide : le spectre d'un anneau strictement hensélien n'a aucune
cohomologie supérieure, et les conséquences que la page en tire —
$\operatorname{cd}_{\ell}(X) \leqslant n+1$, et
$\operatorname{cd}_{\ell}(Z) \leqslant n+1-\operatorname{codim}(Z,X)$ pour $Z$
fermé — sont alors vraies mais sans contenu. L'argument du Lemme 1, qui
travaille sur un sous-schéma $S \subset X$ et sur $\mathbf{P}^1_Y$, montre que
l'objet visé est un ouvert de $\operatorname{Spec} A$ ou un schéma de type fini
au-dessus de lui ; le feuillet ne permet pas de dire lequel, et la présente
lecture ne le répare pas.}

\textbf{Énoncé $2_n$.} Pour $A$ local noethérien de dimension $n$, $f \in A$ et
$X = \operatorname{Spec} A$, on a $\operatorname{cd}_{\ell}(X_f) \leqslant n$,
pour tout $\ell$ premier à la caractéristique résiduelle.

Puis les trois lemmes qui les relient.

\textbf{Lemme 1.} $2_{n'}$ pour tout $n' \leqslant n+1$ entraîne $1_n$.

La démonstration est un joli tour. Pour $Y \subset X$ fermé de codimension $m$,
on compactifie dans $\bar{X} = \mathbf{P}^1_Y$ ; on pose $T = \bar{S} - S$, qui
est strictement local ; et l'on écrit la suite longue reliant $H^i(\bar{S},F)$,
$H^i(S,F)$ et $H^{i+1}_T(\bar{S},F)$. Le premier terme se calcule sur la fibre
spéciale, $H^i(\mathbf{P}^1_k, F_0)$, nul dès $i \geqslant 3$ ; le troisième
tombe sous l'hypothèse $2_{n'}$, appliquée avec $n' = n+1-m$. Il reste à
vérifier que la dimension de $\bar{S}_x$ ne dépasse pas $n+1-m$, ce que donne
l'inégalité
\[
\dim \bar{S}_x + \operatorname{codim}\bigl(\bar{S}_x, \bar{X}_x\bigr)
\leqslant \dim \bar{X}_x = n+1
\]
et le fait que la codimension y est $\geqslant m$.

\textbf{Lemme 2.} $1'_n$ et $2'_{n+1}$ entraînent $2'_{n+2}$.

On plonge $A$ dans une algèbre de séries $k\langle x_1, \dots, x_{n+2}\rangle$,
on fibre $\operatorname{Spec}$ de celle-ci sur
$\operatorname{Spec} k\langle x_1, \dots, x_{n+1}\rangle$ — une fibration de
dimension relative $1$ — et l'on applique $1'_n$ aux fibres pour conclure que
$R^q f_{*}(F)$ est nul en codimension $< q-1$. La suite spectrale donne alors
\[
H^p\bigl(U, R^q f_{*}(F)\bigr) \neq 0 \Longrightarrow
p \leqslant n+1-(q-1), \quad \text{soit } p+q \leqslant n+2,
\]
d'où l'annulation voulue.

\textbf{Lemme 3.} $1'_n$ entraîne : pour $X$ affine de type fini sur un corps
séparablement clos $k$, de dimension de Krull $d \leqslant n+1$,
\[
\operatorname{cd}_{\ell}(X) \leqslant \dim X.
\]

\emph{C'est le théorème d'annulation d'Artin}, et c'est le sommet du dossier.
La démonstration est celle qui a été annoncée dans le résumé, et elle tient en
quatre lignes une fois les lemmes en place. Récurrence sur $d$ ; le cas
$d \leqslant 0$ est trivial ; on se ramène à $X = \mathbf{A}^d_k$ ; on fibre
\[
\mathbf{A}^d_k \xrightarrow{\ f\ } \mathbf{A}^{d-1}_k
\]
avec des fibres de dimension $1$ ; l'hypothèse $1'_n$ appliquée aux fibres donne
que $R^q f_{*}(F)$ est nul en codimension $\leqslant q-1$ ; et l'hypothèse de
récurrence en dimension $d-1$ donne
\[
H^p\bigl(Y, R^q f_{*}(F)\bigr) \neq 0 \Longrightarrow p \leqslant n-(q-1),
\quad \text{soit } p+q \leqslant n+1 .
\]
Tous les $E_2^{pq}$ hors de ce cône sont nuls, et la conclusion tombe.

Les pages 77 à 80 poursuivent vers une forme relative, sur une base strictement
locale plutôt que sur un corps :

\textbf{Théorème (page 79).} Soient $A$ strictement local, $Y =
\operatorname{Spec} A$, $Y' = Y - \{\text{point fermé}\}$, $X$ affine de type
fini sur $Y$, et $X' = X \times_Y Y'$. Alors
\[
\operatorname{cd}_{\ell}(X) \leqslant \operatorname{Sup}
\bigl[\,2\dim X_0,\ 2\dim X' + 1\,\bigr],
\]
$X_0$ désignant la fibre spéciale.\footnote{La page écrit d'abord
« Question », puis biffe le mot et le remplace par « Théorème » — la trace, dans
l'encre, du moment où il a cessé de douter.} Une note marginale ajoute le cas
qui intéresse : si $X$ est intègre et domine $Y$, on trouve
$\operatorname{cd}(X) \leqslant \dim Y + d$, où $d$ est la dimension de la fibre
générique.

La page 80 ouvre la récurrence sur $r$ qui devrait établir ce dernier point, en
plongeant $X$ dans $\mathbf{A}^r_Y$ et en comparant les codimensions selon que
$X$ rencontre ou non la fibre spéciale. Elle ne la conduit pas à son terme.

\pagerange{81}{84}
\subsection*{Ce qui reste conjectural (pages 81 à 84)}

Trois pages enregistrent l'état de la question une fois les théorèmes acquis :
ce qu'on voudrait, ce qu'on sait faux, et ce qu'on ne sait pas.

\textbf{Ce qu'on voudrait.} Pour $X$ noethérien affine, en posant
\[
n = \operatorname*{Sup}_{x \in X}\bigl[\operatorname{cd}_{\ell}(k(x))
+ \dim \mathcal{O}_{X,x}\bigr],
\]
a-t-on $\operatorname{cd}_{\ell}(X) \leqslant n$ ? C'est la forme la plus
naturelle du théorème d'Artin sur une base quelconque.

\textbf{Ce qu'on sait faux.} La page 81 se clôt sur un contre-exemple — un $X$
pour lequel $n = 2$ alors que $\operatorname{cd}(X) = 3$.\footnote{Les
conditions de l'exemple sont perdues : la ligne qui les porte est raturée et
réécrite, et seul le décompte final se lit. La lecture enregistre que le
dossier tient la majoration naïve pour fausse, sans pouvoir dire sur quel
exemple.}

\textbf{Le cas arithmétique.} La page 82 traite à part les schémas affines sur
$\mathbf{Z}$, qui sont le cas où la base a elle-même de la cohomologie. Pour
$f : X \to S = \operatorname{Spec}(\mathbf{Z})$ avec $\ell$ inversible sur $X$,
on écrit la suite spectrale de Leray et l'on sépare deux majorations : au point
générique de $S$, $R^q f_{*}(F)$ est nul pour $q > d$ ; en un point fermé, pour
$q > 1+d$. Le décompte des trois zones — $q > 1+d$ ; $q = 1+d$, où le faisceau
est concentré sur les points fermés et donc $p \leqslant 1$ ; $q \leqslant d$,
où $p \leqslant 2$ — donne
\[
\operatorname{cd}_{\ell}(X) \leqslant d + 2 = \dim X + 1 .
\]
Le $+1$ est le prix de l'arithmétique : $\operatorname{Spec}(\mathbf{Z})$ se
comporte comme une courbe de plus.

\textbf{Les conjectures de la page 84.} Sous le titre « Conjectures sur la
cohomologie des schémas affines », pour $X$ affine intègre et $\ell$ premier :
\begin{enumerate}
\item[$1^{\circ}$)] si $\ell$ est inversible sur $X$,
$\operatorname{cd}_{\ell}(X) \leqslant \operatorname{cd}_{\ell}(k(x))$, avec en
marge la question de savoir si c'est une égalité ;
\item[$2^{\circ}$)] si $\ell$ n'est ni inversible ni nilpotent,
\[
\operatorname{cd}_{\ell}(X) \leqslant
\operatorname{Sup}\bigl(\operatorname{cd}_{\ell}(k(x)),\ \nu(V(\ell)) + 3\bigr),
\]
où, pour un schéma $Z$ de caractéristique $\ell$, $\nu(Z)$ est le supremum des
$\nu(k(z))$ aux points maximaux, et où $\nu(k)$ est le \emph{degré
d'imperfection} de $k$, défini par $\ell^{\nu(k)} = [k : k^{\ell}]$ ;
\item[$4^{\circ}$)] pour $A$ local hensélien intègre de caractéristique
résiduelle $\ell$ et de corps des fractions de caractéristique $0$, et $U$
ouvert dans $\operatorname{Spec} A$,
\[
\operatorname{cd}_{\ell}(U) \leqslant \dim A + \nu(k)
+ \operatorname{cd}_{\ell}(k).
\]
\end{enumerate}

L'invariant $\nu$ mérite d'être remarqué : en caractéristique $\ell$, ce n'est
plus la dimension qui gouverne mais le défaut de perfection du corps de base, et
c'est lui qui remplace le degré de transcendance dans les majorations. Une
conjecture intermédiaire, numérotée 1 bis), est écrite trois fois sur la page,
chaque version raturant la précédente ; elle n'est pas arrêtée.

Une note de sa main, portée sur un feuillet voisin, donne le seul cas où la
difficulté se dissipe : \emph{pour les bons anneaux locaux henséliens de
caractéristique résiduelle nulle, tout se simplifie grâce à Artin} — c'est-à-dire
grâce au théorème d'approximation.

\pagerange{89}{90}
\section*{Deux conséquences géométriques}

\subsection*{Lefschetz faible, qui consomme l'annulation affine (pages 89 et 90)}

Deux pages écrites au net montrent à quoi sert le théorème qu'on vient
d'établir, et c'est l'endroit du dossier où la mécanique se voit le mieux.

Soient $X$ projectif régulier sur un corps séparablement clos $k$, de dimension
$n$ ; $Y$ une section hyperplane régulière ; $F$ un faisceau de $\ell$-torsion
localement constant sur $X$, avec $\ell$ premier à la caractéristique de $k$.
Posons $U = X - Y$ — qui est \emph{affine}, puisque complémentaire d'une section
hyperplane. La suite exacte de cohomologie à support, avec l'isomorphisme de
Gysin $H^i_Y(X,F) \simeq H^{i-2}(Y, F_Y \otimes \check{T}_{\ell})$, s'écrit :

\begin{tikzcd}[column sep=small, row sep=small, nodes={font=\scriptsize}]
H^{i-1}(U,F) \arrow[r] & H^{i}_{Y}(X,F) \arrow[r] \arrow[d, no head, "\wr"] & H^{i}(X,F) \arrow[r] & H^{i}(U,F) \arrow[r] & H^{i+1}_{Y}(X,F) \arrow[d, no head, "\wr"] \\
 & H^{i-2}(Y, F_{Y} \otimes \check{T}_{\ell}) & & & H^{i-1}(Y, F_{Y} \otimes \check{T}_{\ell})
\end{tikzcd}

Et l'on y injecte le théorème d'annulation : $H^i(U,F) = 0$ pour $i > n$,
puisque $U$ est affine de dimension $n$. Deux termes de la suite disparaissent
donc au-delà du rang $n$, et il reste :

\textbf{Théorème de Lefschetz.} L'application
$H^{i-2}(Y, F_Y \otimes \check{T}_{\ell}) \to H^{i}(X,F)$ est surjective pour
$i \geqslant n+1$ et bijective pour $i \geqslant n+2$.

\textbf{Corollaire 1.} Si $X$ est connexe régulier de dimension $n$,
$H^{2n}(X, \mu_N^{\otimes n}) \simeq \mathbf{Z}/N\mathbf{Z}$ pour tout $N$
premier à la caractéristique. La démonstration est une récurrence sur $n$ :
comme $2n \geqslant n+2$ dès $n \geqslant 2$, le théorème s'applique et donne
\[
H^{2n}(X, \mu_N^{\otimes n}) \simeq
H^{2n-2}(Y, \mu_N^{\otimes n} \otimes \check{T}) =
H^{2n-2}(Y, \mu_N^{\otimes (n-1)}),
\]
le twist et son dual se compensant d'un cran.

\textbf{Corollaire 2.} L'application de restriction $H^j(X,F) \to H^j(Y,F)$ est
injective pour $j \leqslant n-1$ et bijective pour $j \leqslant n-2$.

C'est la forme sous laquelle le théorème de Lefschetz faible est cité
aujourd'hui, et l'on voit ici qu'il n'est pas un théorème indépendant : c'est
l'annulation sur l'ouvert affine $U$, lue à travers la suite de Gysin. Le
dossier ne le dit pas en ces termes, mais son organisation le montre — c'est
pourquoi ces deux pages suivent immédiatement les conjectures affines.

\pagerange{92}{98}
\subsection*{Spécialisation et asphéricité locale (pages 92 à 98)}

La dernière suite du dossier, annoncée par une chemise de sa main portant
« Théorème de Spécialisation : Corollaire », vise un résultat d'une autre
nature : non plus une borne sur la cohomologie, mais un énoncé disant qu'un
morphisme lisse ne crée pas de revêtements. Les feuillets sont écrits très vite
et leur prose est en grande partie perdue ; les énoncés et les formules passent,
et c'est d'eux que la lecture rend compte.

\textbf{Lemme 1 (une pureté).} Soient $A$ local noethérien, quotient d'un anneau
régulier, $t \in \mathfrak{m}_A$ un élément $A$-régulier tel que le complété
$t$-adique soit normal, et $\dim A \geqslant 3$. Soit $B$ une extension finie
normale de $A$ telle que $B/tB$ soit étale sur $A/tA$ hors de l'origine. Alors
$B$ est étale sur $A$ hors de l'origine.

La démonstration passe par $X' = \operatorname{Spec} A - \{a\}$ et l'ouvert $U$
de $X'$ où $B$ est étale : l'hypothèse donne
$(X' - U) \cap V(t) \subset \{a\}$, donc $\dim(X' - U) \leqslant 1$, et il
s'agit de prolonger un revêtement étale de $U$ à $X'$ tout entier. Ce qui
l'autorise est la condition de profondeur aux points de $X' - U$, où
$\dim \mathcal{O}_{X,x} \geqslant 2$ — c'est le mécanisme de la pureté de
Zariski--Nagata, et une note au crayon en marge le nomme : « condition des
chaînes ».

\textbf{Lemme 2.} Soient $f : X \to Y$ lisse à fibres de dimension $1$, $s$ une
section de $X$ sur $Y$ avec $Y$ intègre et $X$ normal, et $X' \to X$ un
revêtement galoisien de groupe $\Gamma$ d'ordre premier aux caractéristiques
résiduelles de $Y$. On y considère l'image réciproque $s^{-1}(X')$ — la trace du
revêtement sur la section — et le complémentaire $X - f(Z)$ de l'image d'un
fermé.\footnote{L'énoncé complet du Lemme 2 n'est pas récupérable : la page est
un palimpseste dont seules les données et deux objets se lisent. La lecture
donne ce que le feuillet porte.}

\textbf{Lemme 3 (Abhyankar).} Soient $A$ régulier local de dimension $2$,
$(x,y)$ un système régulier de paramètres, et $B$ une extension normale de $A$
qui est (i) galoisienne d'ordre premier à la caractéristique résiduelle, (ii)
d'indice de ramification constant le long de $V(x)$, (iii) non ramifiée en
dehors. La page l'accompagne d'un croquis — deux diamètres $V(x)$ et $V(y)$ se
coupant au point fermé — et la page 96 donne la conclusion sous forme
explicite : avec $\Gamma = \mathbf{Z}/n\mathbf{Z}$ et $B$ intègre,
\[
B \simeq A[T]/(T^n - x), \qquad
B \otimes_A A' = A'[T]/(T^n - x') \quad \text{pour } A' = A/yA,
\]
$x'$ étant l'image de $x$. C'est le lemme d'Abhyankar, que la page nomme
elle-même : un revêtement modérément ramifié le long d'un diviseur à
croisements normaux est, localement, une extension de Kummer.

Ces trois lemmes en place, la page 96 énonce le théorème :

\textbf{Théorème.} Un morphisme lisse $f : X \to Y$ est localement
$1$-asphérique pour les groupes finis d'ordre premier aux caractéristiques
résiduelles de $Y$.

La démonstration se ramène à $Y$ noethérien, puis, par fibration, au cas de la
dimension relative $1$ — où les trois lemmes s'appliquent. C'est le noyau de ce
qu'on appelle aujourd'hui l'\emph{acyclicité locale des morphismes lisses}, dont
le théorème de changement de base par un morphisme lisse est la forme
cohomologique ; le titre de la chemise, « Théorème de Spécialisation », nomme la
même chose du côté des fibres.

Les applications de la page 97 sont trois, et la troisième explique pourquoi le
résultat était voulu : la constance locale des $R^i f_{*}(F)$ pour $f$ lisse
propre ; l'invariance $H^i(X,F) \xrightarrow{\ \sim\ } H^i(X',F')$ par extension
séparable du corps de base ; et le cas d'un $Y$ lisse dans un $X$ lisse sur une
base.

La page 98 se clôt sur un Corollaire qui met en regard trois familles de
faisceaux — $R^i h_{!}(F'|V)$, $R^i h_{*}(F|V)$ et $R^i g_{*}(F'|Y)$, où
$Y \subset X$ est fermé, $V = X - Y$ son complémentaire ouvert, et
$F' = \mathcal{H}om(F, I_{X/S})$ — et demande leur constance locale
simultanée.\footnote{Les conditions sont numérotées (i), (ii), (iii) sur la
page, dans cet ordre-là : (ii) d'abord, puis (i) reliée à elle par une
accolade, puis (iii). La prose qui devait dire ce qu'elles ont d'équivalent est
perdue.}

\pagerange{11}{17}
\section*{Les feuillets de tapuscrit corrigés}

Onze feuillets du dossier ne sont pas de sa main : ce sont des épreuves
dactylographiées d'autres exposés, qu'il a corrigées à la plume et dont il a
retourné les dos pour y écrire ce qui précède. Ils ne portent pas sur la
dimension cohomologique et ne sont donnés ici ni pour leur contenu propre — on
le trouvera rédigé dans SGA — ni comme pièces d'archive, mais parce que les
corrections sont de lui et que les corrections sont des mathématiques.

\subsection*{Groupes de type multiplicatif : l'isotrivialité (pages 11 à 17)}

Quatre feuillets numérotés 27 à 30 par le dactylographe et rangés à rebours dans
la chemise, portant les énoncés 5.9 à 5.16 d'un exposé sur les groupes de type
multiplicatif. Ils courent vers un seul théorème :

\textbf{Théorème 5.16.} Sur un schéma localement noethérien \emph{normal}, tout
groupe de type multiplicatif de type fini est \emph{isotrivial} — c'est-à-dire
devient diagonalisable après un revêtement fini étale surjectif.

Le chemin passe par l'anti-équivalence entre groupes de type multiplicatif de
type fini et groupes constants tordus à engendrement fini (Corollaire 5.9), par
la représentabilité de $\underline{\operatorname{Isom}}$ par un ouvert-fermé de
$\underline{\operatorname{Hom}}$ (Corollaire 5.10), et par un critère
d'isotrivialité en trois conditions (Proposition 5.11) dont l'équivalence
complète demande que la base soit localement noethérienne. Deux lemmes
techniques portent l'ensemble : tout sous-schéma fermé quasi-compact d'un
schéma constant tordu est fini (Lemme 5.12) ; et sur une base localement
noethérienne connexe, une partie ouverte et fermée d'un schéma constant tordu
quasi-isotrivial dont \emph{une} fibre est finie est finie tout entière (Lemme
5.13). La Remarque 5.15 note que l'hypothèse de normalité peut être affaiblie
en « géométriquement unibranche ».

Les corrections de sa main sont nombreuses et de deux sortes. Les unes sont des
renumérotations en chaîne — le chiffre des unités repassé à l'encre, 5.8 devenu
5.9, 5.9 devenu 5.10, sur les quatre feuillets — avec les références mises à
jour. Les autres sont mathématiques : « somme disjointe de » devient « réunion
de » dans la condition (ii) de 5.11, ce qui n'est pas la même chose ;
« quasi-compacts » remplace la condition de type fini dans le Lemme 5.12 ; et
l'énoncé du Lemme 5.13 voit ses $R$ changés en $P$ et $Z$ pour distinguer le
schéma de la partie ouverte et fermée qu'on y considère — distinction que le
tapuscrit confondait et sans laquelle l'énoncé n'a pas de sens.\footnote{La
démonstration de 5.14, sur le même feuillet, garde le $R$ du tapuscrit là où
l'énoncé porte $P$ corrigé : la correction n'a pas été propagée. La
transcription le note ; la lecture le signale et n'harmonise pas.}

\pagerange{25}{27}
\subsection*{Complétion formelle, et les schémas de Picard non propres (pages 25 et 27)}

Deux feuillets d'un tapuscrit sur la complétion formelle le long d'un fermé,
corrigés de sa main — les accents circonflexes des complétés et les flèches que
la machine ne portait pas sont ajoutés à l'encre.

Le Théorème 2.1 donne un critère de cohérence et de Mittag-Leffler : si le
$\mathcal{S}$-module gradué $\mathcal{K}^i = R^i f_{*}(\operatorname{gr}_
{\mathcal{J}}(\mathcal{F}))$ est de type fini pour $i = n-1, n, n+1$, alors
$R^n f_{*}(F)$ est cohérent, l'homomorphisme de comparaison est un isomorphisme
topologique, et le système projectif des $R^n f_{*}(F_k)$ vérifie la condition
de Mittag-Leffler uniforme. Le Théorème 2.2 traite le cas d'un anneau adique
noethérien dont l'idéal de définition est principal.

Le plus intéressant est la Remarque 5.11, où il annonce un programme. Sous des
hypothèses de pureté et de parafactorialité sur les anneaux locaux des fibres
géométriques, les foncteurs de complétion devraient être des équivalences — sur
les revêtements étales, sur les modules inversibles. Et il ajoute :

\begin{quote}
Utilisant des résultats récents de \textsc{Murre}, il est probable que l'on doit
pouvoir en déduire des théorèmes d'existence des schémas de Picard pour certains
schémas algébriques \emph{non propres}. De façon générale, l'élimination
d'hypothèses de propreté dans divers théorèmes d'existence, notamment de
représentabilité de foncteurs comme les foncteurs de Hilbert, ou de Picard etc,
à l'aide
\end{quote}

Le feuillet s'arrête là, au milieu de la phrase, et la suite n'est pas dans le
dossier.

\pagerange{71}{71}
\subsection*{Quotients par un groupe de type multiplicatif (page 71)}

Un feuillet isolé, du même exposé que les précédents, sur l'existence du
quotient $Y = X/G$ pour $G$ de type multiplicatif : il est affine sur la base,
et hérite des conditions de finitude de $X$. Le Corollaire 2.4 en tire que le
morphisme graphe $X \times_S G \to X \times_S X$ est une immersion fermée, et
qu'il en va de même de $g \mapsto s \cdot g$ pour toute section $s$. Le
Corollaire 2.5 conclut : un monomorphisme $u : G \to H$ de schémas en groupes,
avec $G$ de type multiplicatif et $H$ affine, est une immersion fermée, le
quotient $H/G$ existe et est affine, et $H$ est un fibré principal homogène sur
lui.

\pagerange{83}{87}
\subsection*{Cohomologie de Hochschild et densité schématique (pages 83 à 87)}

Trois feuillets du même tapuscrit, barrés de longues diagonales, dont les
lettres de groupe et plusieurs indices sont encrés à la main là où la machine
avait laissé des blancs.

\textbf{Théorème 3.1.} Pour $S$ affine, $G$ un $S$-groupe de type multiplicatif
et $F$ un faisceau quasi-cohérent sur lequel $G$ opère, la cohomologie de
Hochschild s'annule : $H^i(G,F) = 0$ pour $i > 0$.

La démonstration est un modèle de descente. La cohomologie de Hochschild se
calcule par un complexe explicite de $A$-modules dont la formation commute au
changement de base ; pour $A'$ plat sur $A$, on a donc
$H^i(G',F') = H^i(G,F) \otimes_A A'$ ; pour $A'$ fidèlement plat, l'annulation
descend ; et cela ramène au cas diagonalisable, qui est acquis.

\textbf{Théorème 3.2.} Sur une base affine, deux homomorphismes $u, v : H \to G$
vers un groupe de type multiplicatif qui coïncident modulo un idéal de carré nul
sont conjugués : il existe $g \in G(S)$ avec $v = \operatorname{int}(g) \circ u$
et $g_0 = e$. C'est un énoncé de relèvement infinitésimal, et il découle de 3.1.

\textbf{Corollaire 4.6.} Soit $X$ un $S$-schéma plat et $U$ un ouvert de $X$ tel
que $U_s$ soit schématiquement dense dans $X_s$ pour tout $s \in S$. Si $X$ est
localement noethérien, ou bien si $X$ est localement de présentation finie sur
$S$ avec $U$ rétrocompact dans $X$, alors $U$ est schématiquement dense dans $X$.

La démonstration est une belle application du « procédé breveté » de descente
noethérienne : on écrit $A$ comme limite inductive de ses sous-$\mathbf{Z}$-
algèbres de type fini, on descend $X$ et $U$ à un niveau $i$, et l'on observe
que l'ensemble $E_i$ des points où la densité schématique a lieu est
\emph{constructible} — parce que la condition revient à ce que $U_{i,s}$ contienne
les points associés du faisceau structural de $X_{i,s}$. Comme la condition est
invariante par extension du corps de base, l'image réciproque de $E_i$ est tout
$S$, donc $E_j$ est tout $S_j$ pour un $j$ assez grand.

Les pages 87 portent trois énoncés supplémentaires numérotés 4.7 et 4.9, plus
une Remarque 4.9 ; les corps en sont en grande partie effacés sous les
diagonales d'annulation.\footnote{Les deux corollaires portent le même numéro
4.9 sur le feuillet, et le chiffre du premier (4.7) est réencré par-dessus la
frappe. La transcription les donne tels quels ; la lecture ne les régularise pas
davantage.}

\pagerange{60}{60}
\section*{Deux fragments}

\subsection*{Normes et degré d'imperfection (page 60)}

Un feuillet au crayon, barré de deux traits, sans rapport apparent avec ce qui
l'entoure. Il porte une condition sur une extension de corps $k \to K$ : pour
que l'invariant $\delta(L/K, k/\ell)$ s'annule pour toute extension $L$ de $K$,
il faut et il suffit que
\[
k \subset \ell(K^p), \quad \text{soit } k(K^p) \subset \ell(K^p),
\]
et l'exemple qui montre la nécessité — $L = K(x)$ avec $x^p \in K$ et
$x^p \notin \ell(K^p)$ — accompagné de l'application
$\Omega^1_{K/k} \otimes_K L \to \Omega^1_{L/k}$ qui mesure le défaut. C'est le
même cercle d'idées que l'invariant $\nu$ de la page 84 : en caractéristique
$p$, ce qui compte n'est pas la dimension mais le défaut de perfection.

\pagerange{100}{100}
\subsection*{Une suite spectrale (page 100)}

Une note de sa main sur le dos d'un feuillet dactylographié étranger, qui pose,
pour une immersion ouverte $i : V \to X$ de complémentaire fermé et $g$ le
morphisme structural, la suite spectrale
\[
R^p g_{*}\bigl(R^q i_{*}(G)\bigr) \Longrightarrow R^{*} h_{*}(G)
\]
reliant les images directes supérieures le long de $h : V \to S$ à celles prises
en deux temps.\footnote{Les deux exposants sont écrits $i$ tous les deux sur le
feuillet ; la lecture les distingue en $p$ et $q$, ce que le sens exige. Le
reste de la note n'est pas lisible.} C'est l'outil dont le Corollaire de la
page 98 a besoin, et il est probable que la note s'y rapporte.

\end{document}
